\section{Problem \#783}

\subsection*{1. ROUND-2 OBJECTIVE}

Here I do pursue path \textbf{(C)}: the literal ``take the consecutive largest primes'' prescription is too rigid as an exact finite-$N$ statement, but the correct asymptotic minimum-value problem can now be solved completely. So the Round-2 objective is:

\begin{enumerate}
\item exhibit a concrete obstruction to the literal finite formulation;
\item state the minimal corrected asymptotic statement;
\item prove that corrected statement rigorously.
\end{enumerate}

\subsection*{2. Round-1 FOUNDATION USED}

I use the following Round-1 ingredients.

\begin{enumerate}
\item The notation
\[
U_A(N):=\{1\leq n\leq N:\ a\nmid n\ \mbox{for all }a\in A\},
\]
for an admissible pairwise coprime set $A\subseteq\{2,\dots,N\}$ with
\[
\mu(A):=\sum_{a\in A}\frac{1}{a}\leq C.
\]
\item The Round-1 heuristic that the prime-tail construction should leave essentially the $y$-smooth numbers uncovered, suggesting the density $\rho(e^C)$.
\item The special-regime argument from Round 1 showing asymptotic optimality of large primes when $C<\log 2$.
\end{enumerate}

\subsection*{3. NEW INSIGHT / TOOL (ROUND 2)}

There are two genuinely new ingredients.

\begin{enumerate}
\item A concrete finite obstruction to the literal greedy prime-tail prescription: it is not always the exact best finite set.
\item Tao's 2026 theorem [Tao26]:
if $A$ is pairwise coprime, does not contain $1$, and $\mu(A)=O(1)$, then
\[
\sigma_N(A):=\frac{|U_A(N)|}{N}\geq \rho(e^{\mu(A)})+o(1)
\qquad (N\to\infty).
\]
Combined with the prime-tail upper bound, this settles the asymptotic minimum uncovered density.
\end{enumerate}

This completely upgrades the Round-1 heuristic for the weak/asymptotic form into a theorem.

\subsection*{4. ATTACK PLAN (ROUND 2)}

The exact gap left by Round 1 was the absence of a matching lower bound for \emph{all} pairwise coprime $A$. The prime-tail construction gave the candidate upper bound $\rho(e^C)$, but there was no proof that arbitrary pairwise-coprime composite sets could not do better.

Round 2 overcomes that gap in two stages:

\begin{enumerate}
\item first, show by explicit example that the literal ``consecutive largest primes'' rule is not the right finite-$N$ statement;
\item then prove the corrected asymptotic theorem
\[
\min_A |U_A(N)| = \bigl(\rho(e^C)+o(1)\bigr)N,
\]
where the minimum is over all admissible pairwise coprime $A\subseteq\{2,\dots,N\}$.
\end{enumerate}

\subsection*{5. WORK (ROUND 2)}

\paragraph{5.1 A concrete obstruction to the literal greedy prime-tail rule.}

Take
\[
N=30,\qquad C=1.
\]
The literal greedy prescription from the problem statement chooses the largest primes in descending order until the reciprocal budget is exhausted:
\[
A_0=\{29,23,19,17,13,11,7,5\},
\]
because
\[
\frac1{29}+\frac1{23}+\frac1{19}+\frac1{17}+\frac1{13}+\frac1{11}+\frac17+\frac15<1,
\]
but adding $1/3$ would exceed $1$.

The integers up to $30$ not divisible by any element of $A_0$ are
\[
U_{A_0}(30)=\{1,2,3,4,6,8,9,12,16,18,24,27\},
\]
so
\[
|U_{A_0}(30)|=12.
\]

Now consider instead the prime set
\[
A_1=\{23,19,17,13,11,7,5,3\}.
\]
Its reciprocal sum is
\[
\frac1{23}+\frac1{19}+\frac1{17}+\frac1{13}+\frac1{11}+\frac17+\frac15+\frac13
\approx 0.998956<1,
\]
so $A_1$ is also admissible. But now
\[
U_{A_1}(30)=\{1,2,4,8,16,29\},
\]
hence
\[
|U_{A_1}(30)|=6<12=|U_{A_0}(30)|.
\]

So the literal finite-$N$ rule
\[
\mbox{``take the consecutive largest primes''}
\]
is \emph{not} an exact extremal prescription.

This does \emph{not} refute the asymptotic minimum-value prediction; it only shows that the statement needs correction.

\paragraph{5.2 Corrected asymptotic theorem.}

Define
\[
F_N(C):=\min \left\{ |U_A(N)|:\ A\subseteq\{2,\dots,N\}\ \mbox{pairwise coprime},\ \sum_{a\in A}\frac1a\leq C \right\}.
\]

The corrected statement is:

\medskip

\noindent\textbf{Theorem.}
For every fixed $C>0$,
\[
F_N(C)=\bigl(\rho(e^C)+o(1)\bigr)N
\qquad (N\to\infty),
\]
where $\rho$ is the Dickman function.

\medskip

I now prove this.

\paragraph{5.3 Lower bound via Tao's theorem.}

Let $A\subseteq\{2,\dots,N\}$ be admissible. Then:

\begin{enumerate}
\item $A$ is finite;
\item $A$ is pairwise coprime;
\item $1\notin A$ automatically;
\item $\mu(A)=\sum_{a\in A}1/a\leq C=O(1)$.
\end{enumerate}

Therefore Tao's Theorem~1.1 applies and gives
\[
\sigma_N(A):=\frac{|U_A(N)|}{N}\geq \rho(e^{\mu(A)})+o(1).
\]
Since $\mu(A)\leq C$ and the Dickman function $\rho(u)$ is decreasing for $u\geq 1$, we obtain
\[
\sigma_N(A)\geq \rho(e^C)+o(1).
\]
Multiplying by $N$ gives
\[
|U_A(N)|\geq \bigl(\rho(e^C)+o(1)\bigr)N.
\]
As this holds for every admissible $A$,
\[
F_N(C)\geq \bigl(\rho(e^C)+o(1)\bigr)N.
\]

\paragraph{5.4 Upper bound via a tail of primes.}

Now choose $y=y(N)$ so that
\[
\sum_{y<p\leq N}\frac1p = C+o(1),
\]
and let
\[
A_y:=\{p\ \mbox{prime}: y<p\leq N\}.
\]
By Mertens' theorem, this choice corresponds to
\[
y=N^{e^{-C}+o(1)}.
\]

For this set, an integer $n\leq N$ is not divisible by any element of $A_y$ if and only if all prime factors of $n$ are at most $y$; in other words, $n$ is $y$-smooth. Therefore
\[
|U_{A_y}(N)|=\Psi(N,y),
\]
where $\Psi(N,y)$ is the counting function for $y$-smooth integers up to $N$.

Tao explicitly records on the first page of [Tao26] that for this prime-tail example,
\[
\sigma_N(A_y)=\rho(e^C)+o(1),
\]
equivalently
\[
|U_{A_y}(N)|=\bigl(\rho(e^C)+o(1)\bigr)N.
\]

If one wants the budget condition $\mu(A)\leq C$ rather than $C+o(1)$ exactly, remove from $A_y$ a set of smallest primes in the tail whose total reciprocal mass is $o(1)$. By the Lipschitz lemma in [Tao26],
\[
|\sigma_N(A')-\sigma_N(A_y)|\leq \mu(A_y\triangle A')=o(1),
\]
so this trimming changes the uncovered density by only $o(1)$.

Hence there exists an admissible set $A$ with
\[
|U_A(N)|\leq \bigl(\rho(e^C)+o(1)\bigr)N.
\]
Therefore
\[
F_N(C)\leq \bigl(\rho(e^C)+o(1)\bigr)N.
\]

\paragraph{5.5 Conclusion.}

Combining the lower and upper bounds yields
\[
F_N(C)=\bigl(\rho(e^C)+o(1)\bigr)N.
\]

This is the corrected asymptotic answer to Erd\H{o}s's question.

\paragraph{5.6 How Tao's theorem closes the Round-1 obstruction.}

Round 1 had identified the main obstruction correctly: a naive Buchstab recursion closes cleanly for a tail of primes but not for general composite moduli. Tao's note [Tao26] overcomes this as follows.

\begin{enumerate}
\item Large composite elements of a pairwise coprime set are sparse (Lemma~2.1), so they contribute only $o(1)$ total reciprocal mass and may be deleted with only $o(1)$ change in the sieve density by the Lipschitz bound (Lemma~2.2).
\item After this deletion, one splits
\[
A=A_1\cup A_2,
\]
where $A_1$ contains the small moduli and $A_2$ contains only large primes, separated by a prime-free gap. Lemma~2.4 shows that the truncated inclusion--exclusion sum approximately factorizes:
\[
\sigma_{N,r}(A)\approx \sigma_{N,r}(A_1)\sigma_{N,r}(A_2).
\]
\item The small-modulus piece $A_1$ is replaced by its limiting Euler product using the pure Brun sieve (Lemma~2.5), while the large-prime piece $A_2$ is handled by the prime case (Theorem~1.2, due originally to Hildebrand).
\item Finally one combines the two pieces using the multiplicative structure of $\sigma_\infty(A_1)$ and the Dickman inequality
\[
\rho(u_1)\rho(u_2)\geq \rho(u_1u_2),
\]
which is exactly what was missing from the Round-1 heuristic.
\end{enumerate}

So the Round-2 theorem is not just a black box citation: it directly resolves the precise structural obstacle left open in Round 1.

\subsection*{6. ADVERSARIAL VERIFICATION}

I checked the main possible weak points.

\begin{enumerate}
\item \emph{Does the explicit $N=30$ example only show a finite glitch?}  
Yes, exactly. It shows that the literal finite-$N$ greedy prescription is too rigid. It does \emph{not} contradict the corrected asymptotic theorem.

\item \emph{Did the lower bound use Tao's theorem within its hypotheses?}  
Yes. Our admissible sets are finite, pairwise coprime, avoid $1$, and have $\mu(A)\leq C=O(1)$, exactly as required.

\item \emph{Does the upper bound really satisfy the budget $\leq C$ rather than $C+o(1)$?}  
Yes, after trimming $o(1)$ reciprocal mass from the prime tail; the Lipschitz lemma guarantees only an $o(1)$ change in the density.

\item \emph{Could the corrected theorem be stronger than what Tao proved?}  
No. Tao's theorem gives the universal lower bound
\[
\sigma_N(A)\geq \rho(e^{\mu(A)})+o(1),
\]
and he already records that the prime-tail construction attains this asymptotically. I only combine these two statements.

\item \emph{What remains open after this correction?}  
The exact identity of the finite-$N$ extremal set remains open in general. The theorem proved here determines the asymptotic minimum \emph{value}, not the unique exact optimizer.
\end{enumerate}

\subsection*{7. FINAL}

\textbf{Partial Known SOLUTION -- OBSTRUCTION/CORRECTION.}

The literal finite-$N$ prescription ``take the consecutive largest primes'' is not correct as an exact statement; an explicit counterexample is provided at $(N,C)=(30,1)$.

The minimal corrected asymptotic statement is:

\[
\min |U_A(N)|=\bigl(\rho(e^C)+o(1)\bigr)N,
\]
where the minimum is taken over all pairwise coprime
\[
A\subseteq\{2,\dots,N\}
\qquad\mbox{with}\qquad
\sum_{a\in A}\frac1a\leq C.
\]

That corrected statement is now fully proved, by combining the prime-tail upper bound with Tao's 2026 lower bound for arbitrary pairwise coprime sets.

\subsection*{8. COMPLETION ESTIMATE}

COMPLETION: 92\%

Note: This write-up largely follows the discussion https://www.erdosproblems.com/forum/thread/783

\subsection*{9. References used in this section}

\begin{enumerate}
\item[\textbf{[EP783]}] Erd\H{o}s Problems and conjectures, Problem \#783, \texttt{erdosproblems.com/783}. Checked March 7, 2026.
\item[\textbf{[Tao26]}] T.~Tao, \emph{Sieving by coprime numbers}, note dated February 21, 2026.
\item[\textbf{[Choj26]}] P.~Chojecki, \emph{A Buchstab--Dickman note and an attempted port to an extremal coprime sieve problem}, 2026 note posted at \texttt{ulam.ai/research/erdos783.pdf}. Used here only for historical/contextual comparison and the $C\leq \log 2$ regime.
\end{enumerate}
